\section{Betrugserkennung bei Kfz-Schadensfällen}
\label{sec:fraud}

Beim ersten spezifischen Data Science Use Case, 
der durch zu entwickelnde Datenarchitektur abgedeckt werden soll,
handelt es sich um die KI-gestütze Betrugserkennung bei Kfz-Schadensfällen.
Versicherungsbetrug verursacht nicht nur erhebliche finanzielle Schäden,
sondern wird durch die rasanten Fortschritte generativer KI 
gleichzeitig einfacher und schwieriger erkennbar.
Gefälschte Schadensfotos und andere Nachweise werden zunehmend hochwertiger.
Und durch die breite Verfügbarkeit performance-starker KI Modelle 
braucht es außerdem keine besonderen Fertigkeiten mehr, 
um Fotos und Videos zu manipulieren \citep{hever2025}.
Um die finanziellen Verluste durch Betrug bei Kfz-Schadensfällen zu minimieren,
soll für die Allianz SE ein KI-Modell entwickelt werden,
welches alle für den Schadenfall relevanten Informationen verarbeitet,
und basierend darauf eine Einschätzung gibt, 
ob es sich dabei um einen Betrugsversuch handelt.
Die für diesen Use Case relevanten Analysedaten sind in 
Tabelle~\ref{tab:use-case-1-analysedaten} fachlich aufbereitet.

\begin{table}[H]
  \centering
  \caption{Analysedaten für die Betrugserkennung bei Kfz-Schadensfälen}
  \label{tab:use-case-1-analysedaten}
  \begin{tabularx}{\textwidth}{@{}p{6.5cm}X@{}}
    \toprule
    Datenquelle & Beschreibung \\
    \midrule
    \midrule
    bestätigte Betrugsfälle, historische Schadensdaten, ggf. externe Fraud-Datenbank
      & Trainingsgrundlage für das KI-Modell \\
    \midrule
    Schadensmeldung
      & Ort, Datum, Uhrzeit, Vorgangsbeschreibung \\
    \midrule
    Schadensfotos/-dokumente
      & relevante Anhänge inkl. EXIF-Metadaten \\
    \midrule
    Telemetriedaten (sofern vorhanden)
      & Fahrstil, Bremsverhalten, Fahrzeugnutzung vor dem Unfall, GPS-Daten \\
    \midrule
      Wetter- und Geodaten
      & zum Abgleich mit dem gemeldeten Schadensort (z.B. bei Hagelschäden) \\
    \midrule
      Kaufkraft-/sozioökonomische Daten je Postleitzahl
      & zusätzlicher Kontext für mögliche Mustererkennung in den Betrugsfällen \\
    \bottomrule
  \end{tabularx}
\end{table}

Wichtig ist dabei, dass das Modell ausschließlich entscheidungsunterstützend
eingesetzt wird.
Die Entscheidung, ob eine Schadensmeldung abgelehnt wird, 
soll weiterhin von den Sachbearbeiter:innen im Schadensmanagement 
getroffen werden.
 

